\subsection{舒尔引理}

本节中，希尔伯特空间均为复的。

命题 6（舒尔引理）。--- 设 \(\varrho\) 是拓扑群 G 在非零希尔伯特空间 E 中的酉表示。那么，\(\varrho\) 不可约，当且仅当 \(\mathrm{Hom}_{\mathrm{G}}(\varrho, \varrho)\) 等于 \(\mathbf{C} \cdot 1_{\mathrm{E}}\)。

\(\mathrm{Hom}_{\mathrm{G}}(\varrho, \varrho)\) 是 \(\mathcal{L}(\mathrm{E})\) 的含单位星子代数（V, p. 380 的引理 5）。对于 E 的每个子表示 F，像为 F 的正交投影 \(p\) 是星代数 \(\mathrm{Hom}_{\mathrm{G}}(\varrho, \varrho)\) 中的幂等元（V, p. 383 的命题 4）。若该代数等于 \(\mathbf{C} \cdot 1_{\mathrm{E}}\)，这意味着 \(p = 0\) 或 \(p = 1_{\mathrm{E}}\)，也即 F = 0 或 F = E。因此 \(\varrho\) 不可约。

反之，假设 \(\varrho\) 不可约。设 \(u\) 和 \(v\) 是星代数 \(\mathrm{Hom}_{\mathrm{G}}(\varrho, \varrho)\) 中满足 \(uv = 0\) 的交换元素。设 \(u\) 非零。于是 \(u\) 的核 F 定义 \(\varrho\) 的一个不同于 E 的子表示，故 F 必为 \(\{0\}\)；由于 F 包含 \(v\) 的像，推出 \(v = 0\)。因此由 I, p. 113 的命题 10，有 \(\mathrm{Hom}_{\mathrm{G}}(\varrho, \varrho) = \mathbf{C} \cdot 1_{\mathrm{E}}\)。

推论 1. --- 设 π 是拓扑群 G 在希尔伯特空间 E 中的不可约酉表示。设 u 是 E 上的闭部分算子。假设 u 的定义域稠密、dom(u) 在 π 下稳定，并且对所有 g ∈ G 有 u ∘ π(g) = π(g) ∘ u。那么 dom(u) = E，且 u 是恒等映射的标量倍数。

部分算子 \(u^{*} \circ u\) 是 E 上的正自伴部分算子（IV, p. 241 的命题 12）；对于 \(v = 1_{\mathrm{E}} + u^{*} \circ u\) 也同样如此，并且后者是单射的，因为 \(-1 \notin \operatorname{Sp}(u^{*} \circ u)\)（IV, p. 248 的命题 17）。有关系式 \(u^{*} \circ \pi(g) = \pi(g) \circ u^{*}\) 对所有 \(g \in G\) 成立（V, p. 381 的引理 6）；从而关系式 \(v \circ \pi(g) = \pi(g) \circ v\) 对所有 \(g \in G\) 成立。由于 \(v\) 是单射，故关系式 \(v^{-1} \circ \pi(g) = \pi(g) \circ v^{-1}\) 对所有 \(g \in G\) 成立。由于 \(v^{-1}\) 属于 \(\mathcal{L}(\mathrm{E})\)，所以由命题 6，存在 \(\lambda \in \mathbf{C}\) 使得 \(v^{-1} = \lambda 1_{\mathrm{E}}\)。必有 \(\lambda \neq 0\)，这意味着 \(\mathrm{E} = \mathrm{Im}(v^{-1}) = \mathrm{dom}(v) \subset \mathrm{dom}(u)\)，因而 \(\mathrm{dom}(u) = \mathrm{E}\)。由于 \(u\) 是闭的，有 \(u \in \mathcal{L}(\mathrm{E})\)，于是 \(u \in \mathrm{Hom}_{\mathrm{G}}(\pi, \pi)\)，并且由命题 6，\(u\) 是一个伸缩变换。

推论 2. --- 设 \(\varrho\) 和 \(\pi\) 是拓扑群 G 在希尔伯特空间 E 和 F 中的不可约酉表示。空间 \(\mathrm{Hom}_{\mathrm{G}}(\varrho, \pi)\) 的维数在 \(\varrho\) 与 \(\pi\) 同构时为 1，否则为零。特别地，若 \(\varrho\) 与 \(\pi\) 同构，则从 \(\varrho\) 到 \(\pi\) 的每个非零 G-态射都是同构。

假设存在非零 G-态射 \(u\)，它从 \(\varrho\) 到 \(\pi\)。线性映射 \(u^{*} \circ u\) 是 \(\mathrm{Hom}_{\mathrm{G}}(\varrho, \varrho)\) 的元素；因此存在复数 \(\lambda\)，使得 \(u^{*} \circ u = \lambda \cdot 1_{\mathrm{E}}\)（命题 6）。

于是有

\[
\langle u (x) \mid u (y) \rangle = \langle x \mid u ^ {*} u (y) \rangle = \lambda \langle x \mid y \rangle
\]

对 E 中所有 \(x\) 和 \(y\) 成立。特别地，\(\lambda \neq 0\)，因为 \(u\) 非零。由于 \(\| u(x)\| = |\lambda |^{1 / 2}\| x\|\) 对所有 \(x\in \mathrm{E}\) 成立，线性映射 \(u\) 是单射，且 \(u\) 的像在 F 中闭（I, p. 107 的引理 8）；于是它是不可约表示 \(\pi\) 的一个非零子表示，从而 \(u\) 是满射。因此，\(u\) 是从 \(\varrho\) 到 \(\pi\) 的同构。于是映射 \(v\mapsto u^{*}\circ v\) 是从空间 \(\mathrm{Hom}_{\mathrm{G}}(\varrho ,\pi)\) 到空间 \(\mathrm{Hom}_{\mathrm{G}}(\varrho ,\varrho)\) 的同构，而后一个空间的维数为 1（命题 6）。

推论 3. --- 设 \(\varrho\) 是拓扑群 G 在希尔伯特空间 E 中的酉表示。表示 \(\varrho\) 不可约，当且仅当空间 \(\mathrm{Sesq}_{\mathrm{G}}(\varrho, \varrho)\) 的维数为 1。此时，该空间由 E 上的内积生成。

根据 V, p. 381 的命题 1，这由命题 6 得出。

推论 4. --- 设 \(\varrho_{1}\) 和 \(\varrho_{2}\) 是拓扑群 G 在希尔伯特空间 E \(_{1}\) 和 E \(_{2}\) 中的非同构不可约酉表示。空间 Sesq \(_{G}\)（\(\varrho_{1}, \varrho_{2}\)）为零。

由 V, p. 381 的命题 1，这由推论 2 得出。

推论 5. --- 设 \(\varrho\) 是拓扑群 G 在希尔伯特空间 E 中的酉表示。设 \(\pi\) 是 G 在希尔伯特空间 F 中的不可约酉表示。

\begin{enumerate}
\def\labelenumi{\alph{enumi})}
\item
  从 \(\pi\) 到 \(\varrho\) 的每个非零 G-态射 u，都存在 \(\lambda \in \mathbf{R}_{+}^{*}\)，使得 \(\lambda u\) 是等距的；
\item
  从 π 到 ρ 的每个 G-态射 u 都有闭像；
\item
  从 \(\varrho\) 到 \(\pi\) 的每个非零 G-态射 v 都是满射。
\end{enumerate}

证明断言 \(a\)，它蕴含断言 \(b\)（I, p. 107 的引理 8）。由于态射 \(u\) 非零，它是单射（V, p. 378 的引理 2）。由公式 \(q(x,y) = \langle u(x)|u(y)\rangle\)，其中 \(x\) 和 \(y\) 属于 F，定义了 F 上的连续半双线性形式。由于 \(u\) 是 G-态射，有 \(q(\pi(g)x,\pi(g)y) = q(x,y)\) 对所有 \(g\)（属于 G）及 \((x,y)\in\mathrm{F}\times\mathrm{F}\) 成立，故 \(q\) 是 G-不变的。由推论 3，存在 \(\alpha\in\mathbf{R}_{+}^{*}\)，使得

\[
q (x, y) = \langle u (x) | u (y) \rangle = \alpha \langle x | y \rangle
\]

对所有 \((x,y)\in\mathrm{F}\times\mathrm{F}\) 成立，因而 \(\alpha^{-1/2}u\) 是等距的。

最后证明断言 \(c\)。由于 \(\pi\) 不可约，\(v\) 的像在 F 中稠密（V, p. 378 的引理 2），故伴随 \(v^*\) 是单射（TVS, V, p. 41 的命题 4）。由 \(b\)，\(v^*\) 的像 H 是 E 的闭子空间；因此它是 \(\varrho\) 的子表示，而由限制到子空间得到的 \(v^*\) 是从 F 到 H 的 G-同构。映射 \(w\) 是 \(v\) 在 H 上的限制，它定义了从 H 到 F 的 G-态射；由于其核是 H 与 \(\text{Ker}(v) = \text{H}^\circ\) 的交，它是单射（同上）。因此 \(w\) 是同构（推论 2）；特别地，\(v\) 是满射。

推论 6. --- 设 \(G_{1}\) 和 \(G_{2}\) 是拓扑群。设 \(\varrho_{1}\) 是 \(G_{1}\) 在希尔伯特空间 \(E_{1}\) 中的不可约酉表示，\(\varrho_{2}\) 是 \(G_{2}\) 在希尔伯特空间 \(E_{2}\) 中的不可约酉表示。外张量积 \(\varrho_{1} \boxtimes \varrho_{2}\) 是 \(G_{1} \times G_{2}\) 在 \(E_{1} \widehat{\otimes}_{2} E_{2}\) 中的不可约酉表示。

由 V, p. 384 的引理 8，外张量积 \(\varrho = \varrho_{1} \otimes \varrho_{2}\) 是 \(G = G_{1} \times G_{2}\) 在空间 \(E = E_{1} \widehat{\otimes}_{2} E_{2}\) 中的酉表示。

设 \(q\) 是 E 上的 \((\mathrm{G}_1 \times \mathrm{G}_2)\)-不变半双线性形式。对于每个 \((x_1, y_1) \in \mathrm{E}_1^2\)，映射 \((x_2, y_2) \mapsto q(x_1 \otimes x_2, y_1 \otimes y_2)\) 属于 \(\operatorname{Sesq}_{\mathrm{G}_2}(\varrho_2, \varrho_2)\)。记 \(b(x_1, y_1)\) 为满足下式的唯一复数：

\[
q (x _ {1} \otimes x _ {2}, y _ {1} \otimes y _ {2}) = b (x _ {1}, y _ {1}) \left\langle x _ {2} \mid y _ {2} \right\rangle
\]

对所有 \((x_{2},y_{2})\in\mathrm{E}_{2}^{2}\) 成立（推论 3）。取 \(\varepsilon\in E_{2}\) 使其范数为 1，于是有 \(b(x_{1},y_{1})=q(x_{1}\otimes\varepsilon,y_{1}\otimes\varepsilon)\) 对所有 \((x_{1},y_{1})\in\mathrm{E}_{1}^{2}\) 成立。该公式说明映射 b 是 \(E_{1}\) 上的半双线性映射。此外

\[
\| b (x _ {1}, y _ {1}) \| \leqslant \| q \| \| x _ {1} \otimes \varepsilon \| \| y _ {1} \otimes \varepsilon \| = \| q \| \| x _ {1} \| \| y _ {1} \|
\]

对所有 \((x_{1},y_{1})\in\mathrm{E}_{1}^{2}\) 成立（EVT, V, p. 26, 公式（5）），因此 b 是连续的。

设 \(g \in G\)，且 \((x_{1}, y_{1}) \in \mathrm{E}_{1}^{2}\)。由于 \(\varrho_{2}\) 是酉的且 q 不变，得到

\[
b \left(\varrho_ {1} (g) x _ {1}, \varrho_ {1} (g) y _ {1}\right) = q \left(x _ {1} \otimes \varrho_ {2} \left(g ^ {- 1}\right) \varepsilon , y _ {1} \otimes \varrho_ {2} \left(g ^ {- 1}\right) \varepsilon\right) = b \left(x _ {1}, y _ {1}\right),
\]

因此 \(b \in \mathrm{Sesq}_{\mathrm{G}_1}(\varrho_1, \varrho_1)\)。于是存在唯一的 \(\lambda \in \mathbf{C}\)，使得 \(b(x_1, y_1) = \lambda \langle x_1 | y_1 \rangle\) 对所有 \((x_1, y_1) \in \mathrm{E}_1^2\) 成立（推论 3），即

\[
q (x _ {1} \otimes x _ {2}, y _ {1} \otimes y _ {2}) = \lambda \langle x _ {1} | y _ {1} \rangle \langle x _ {2} | y _ {2} \rangle
\]

对所有 \((x_{1},y_{1},x_{2},y_{2})\in\mathrm{E}_{1}^{2}\times\mathrm{E}_{2}^{2}\) 成立。由此可推出空间 \(\operatorname{Sesq}_{\mathrm{G}_{1}\times\mathrm{G}_{2}}(\varrho,\varrho)\) 的维数为 1，这意味着表示 \(\varrho=\varrho_{1}\boxtimes\varrho_{2}\) 不可约（同上）。

回忆一下，U 表示模为 1 的复数所成的群。

推论 7. --- 设 \(\pi\) 是拓扑群 G 在希尔伯特空间 E 中的不可约酉表示。存在从 G 的中心 C 到 U 的连续同态 \(\chi\)，满足 \(\pi(z) = \chi(z) \cdot 1_{\mathrm{E}}\) 对所有 \(z \in \mathrm{C}\) 成立。特别地，若 G 是交换群，则 \(\dim(\mathrm{E}) = 1\)。

对于 \(z \in \mathrm{C}\)，映射 \(\pi(z)\) 属于 \(\mathrm{Hom}_{\mathrm{G}}(\pi, \pi)\)；因此它具有 \(\chi(z) \cdot 1_{\mathrm{E}}\) 的形式，其中 \(\chi(z)\) 是某个复数（命题 6）。由于 \(\pi(z)\) 是酉映射，有 \(|\chi(z)| = 1\)。此外，由于 \(\pi\) 是同态，映射 \(z \mapsto \chi(z)\) 是同态。取 E 中 \(v \neq 0\)；从 C 到 U 的映射 \(z \mapsto \chi(z)v = \pi(z)v\) 是连续的，因此同态 \(\chi\) 是连续的。

最后，若 G 是交换群，则 C = G，所以对 G 中所有 g，有 \(\pi(g) = \chi(g) \cdot 1_{\mathrm{E}}\)；于是 E 的每个一维子空间都是 \(\pi\) 的子表示，而由于 \(\pi\) 不可约，空间 E 必为一维。

定义 6. --- 设 π 是拓扑群 G 在希尔伯特空间 E 中的不可约酉表示。满足对 C 中所有 z 有 π(z) = χ(z) · \(1_{\mathrm{E}}\) 的从 G 的中心 C 到 U 的同态 χ，称为 π 的中心特征标。

注. --- 设 \(\pi\)（分别为 \(\varrho\)）是拓扑群 G（分别为 H）在希尔伯特空间 E（分别为 F）中的不可约酉表示。记 \(\chi\)（分别为 \(\eta\)）为 \(\pi\)（分别为 \(\varrho\)）的中心特征标。表示 \(\overline{\pi}\) 的中心特征标是 \(\overline{\chi}\)；而 G × H 的不可约酉表示 \(\pi \boxtimes \varrho\)（推论 6）的中心特征标，是 G × H 上的特征标 \(\chi \boxtimes \eta\)：\((g, h) \mapsto \chi(g)\eta(h)\)。

